\documentclass[../../main.tex]{subfiles}

\graphicspath{{../../image/}{../../figure/}}

\begin{document}

\section{快速振荡与伴随平衡律的稳定化作用}\label{section:ch01-快速振荡}

考虑齐次守恒律方程组
\begin{align}
\operatorname{div} G(U(X))=0\label{align:ch01-homogeneous}
\end{align}
并假设对某状态 $V,W\in\mathcal{O}$ 与 $N\in\mathbb{S}^{k-1}$，有
\begin{align}
[G(W)-G(V)]N=0.\label{align:ch01-振荡跃度}
\end{align}
则可按如下程序在 $\mathbb{R}^{k}$ 上构造\cref{align:ch01-homogeneous} 的高度振荡弱解：从有限族平行、均正交于 $N$ 的 $(k-1)$ 维超平面出发，在相邻超平面之间定义一个在 $\mathbb{R}^{k}$ 上交替取常值 $V$ 与 $W$ 的函数 $U$。显然 $U$ 是\cref{align:ch01-homogeneous} 的弱解，这由\cref{align:ch01-振荡跃度} 与\cref{theorem:ch01-BV解刻画} 可知。

由此可构造一解序列，它在 $L^{\infty}$ 弱$^{*}$ 意义下收敛到形如 $U(X)=\rho(X\cdot N)V+[1-\rho(X\cdot N)]W$ 的某个 $U$，其中 $\rho$ 是任一从 $\mathbb{R}$ 到 $[0,1]$ 的可测函数。显然一般地，除非 $G(\cdot)N$ 恰好沿 $\mathbb{R}^{n}$ 中连接 $V$ 与 $W$ 的直线段为仿射，这样的 $U$ 不再是\cref{align:ch01-homogeneous} 的解。这种不稳定性把可支撑激波阵面的方程组与不能支撑激波阵面的椭圆型方程组区分开来。

现假设 $G$ 配有伴随 $Q$，且 $[Q(W)-Q(V)]N\neq 0$。注意到，施加容许性条件 $\operatorname{div}Q(U)\leqslant 0$ 将使上述构造的振荡解不合格，因为由\cref{theorem:ch01-BV伴随不等式}，它将不允许沿 $N$ 方向同时发生从 $V$ 到 $W$ 与从 $W$ 到 $V$ 的跳跃。因而不等式\cref{align:ch01-admissibility} 似乎起着稳定化作用。这种稳定化在多大程度上有效，将是本书讨论的一个主要问题。

\end{document}
